% =====================================================================
% RAOD++ / iso-cost study — Experimental Setup
% Target: Image and Vision Computing, VSI: VP-NAV
% Every number here is read from output/ ; nothing is typed from memory.
% =====================================================================

\section{Experimental setup}
\label{sec:setup}

\subsection{Datasets}
\label{sec:setup:data}

Two UAV detection benchmarks are used. They share a domain but differ in scale
statistics, which is what makes the second one a test of the mechanism rather than a
repetition of the first.

\textbf{VisDrone2019-DET} is the primary benchmark: 6,471 training and 548 validation
images with 343,197 and 38,759 annotated boxes across ten categories. The class
distribution is heavily skewed, with a ratio of 44.6 between the most and least frequent
training category. Annotations marked as ignored regions and the \emph{others} category
are discarded, boxes are clipped to the image and normalised, and the number of discarded
boxes is recorded. Objects are predominantly small: with the modal image resolution the
median box area falls in the COCO small band.

\textbf{UAVDT} is the second benchmark, used to test whether the findings transfer. It
provides an official split by video sequence, 30 sequences for training and 20 for
evaluation, with no sequence appearing in both. That split is adopted unchanged. Because
consecutive frames of a sequence are near duplicates, a frame-level split would constitute
severe leakage; the sequence-level split avoids it by construction. Frames are then thinned
within each sequence by a fixed stride (5 for training, 9 for evaluation) to bring the set
to a size comparable with VisDrone, yielding 4,840 training and 1,854 evaluation images
with 84,727 and 42,041 boxes over the three official vehicle categories. UAVDT distributes two subsets in one archive, and only one of them is used here. The
sequences named \texttt{M****} are the detection subset, annotated with the three vehicle
categories. The sequences named \texttt{S****} are the tracking subset: one box per frame
carrying a single \emph{vehicle} label, since tracking does not distinguish categories.
Every annotation file in the archive was counted rather than sampled: the archive holds
$846{,}857$ annotated objects, of which $37{,}084$ are the tracking subset's single-class
boxes, one in each of $37{,}084$ frames. Those frames are excluded, because a box without a
category cannot be scored against a three-class detector. The evaluation set is therefore
the $16{,}592$ frames of the detection subset, thinned by stride to $1{,}854$.

Both splits are frozen: the sorted list of file names is hashed with SHA-256 before any
training, and every subsequent run re-computes the hash and aborts on a mismatch. Leakage
audits are run and reported for both datasets, checking for shared file names and, on
UAVDT, for shared sequences. All evaluation in this work uses the frozen validation split
of the corresponding dataset; no result is ever computed on training images.

\subsection{Models and training}
\label{sec:setup:training}

Three architecture variants are used, all from the YOLOv8 family: the nano model
(3.01\,M parameters, 8.20\,GFLOPs at 640\,px), the small model, and the nano model extended
with a P2 detection head, which is the standard remedy for small objects in the UAV
literature. Each is initialised from COCO-pretrained weights, with the detection head
remapped by class name.

Training is identical across conditions apart from the factor under test. Every run uses a
ceiling of 200 epochs with early stopping after 30 epochs without improvement, the default
Ultralytics optimiser and augmentation settings, deterministic mode, and a fixed seed.
The per-image batch is the largest that fits in 24\,GB at the given input resolution and is
recorded per run; gradient accumulation to a nominal batch of 64 keeps the effective batch
constant, so the per-image batch is a memory constraint rather than a treatment.
Input resolution is varied over $\{640, 960, 1280, 1536\}$ for the resolution arm and
$\{640, 960, 1280\}$ for the capacity arm.

The resolution arm's $640$ point on VisDrone is the baseline, which was trained before the
grid script existed and therefore by a different code path. Its recorded settings were
compared field by field against every other condition of the grid
(\texttt{scripts/check\_run\_configs.py}): all experiment-defining parameters --- schedule,
optimiser and its hyperparameters, every augmentation setting, the nominal batch, the
deterministic flag and the loader worker count --- are identical, and the only differences
are the input resolution, the per-image batch and the architecture, which are the factors
under test. The two code paths therefore produce the same configuration.

Three seeds ($42$, $123$, $2024$) are used for every condition of the cost study. The
baseline at $640$\,px on VisDrone was additionally repeated over ten seeds; its
seed-to-seed standard deviation of $0.0015$ mAP$_{50\text{-}95}$ sets the noise scale
against which all differences in this work are judged.

\subsection{Cost measurement}
\label{sec:setup:cost}

The field usually reports parameter counts and FLOPs. Neither predicts wall-clock latency
well in the regime studied here, so cost is measured directly, on the same machine, for
every trained model.

Latency is measured at batch size one, which is how a detector runs on board an aircraft.
Fifty validation frames are decoded into RAM beforehand so that disk access is not timed.
Each measurement performs 100 warm-up inferences followed by 1{,}000 timed inferences with
explicit device synchronisation, and reports the mean and the 95th percentile. The whole
procedure is repeated three times and the median is taken; the spread across the three
passes is stored alongside the value. Measurement is refused unless the GPU is idle, since
figures taken immediately after training, on a hot and fragmented device, were found to
vary by more than 30\% for the same model.

Board power is sampled through NVML during the timed passes, and energy per frame is
derived as the product of mean power and mean latency. This is board power, not system
wall power; the distinction is stated wherever energy is reported.

\subsection{Comparison at equal cost}
\label{sec:setup:isocost}

Comparing named configurations pairwise assumes they sit at comparable cost. Measurement
showed they do not: two conditions intended as a matched pair differed by 18\% in latency,
which silently grants one arm a larger compute budget. Conditions are therefore not
compared pairwise. Instead, each arm --- resolution, capacity, P2 head --- is represented
by its accuracy-versus-cost curve, built per seed from that seed's measured points.

For a pair of arms, the comparison is restricted to the cost interval in which both arms
have real measurements. Within that interval, each arm's curve is interpolated onto a
common grid of 25 budgets, per seed, and the difference is taken seed by seed. No
extrapolation is performed: budgets outside the shared interval are dropped, and if two
arms do not overlap at all they are reported as not comparable. The procedure is applied
independently on the latency axis and on the energy axis.

\subsection{Statistical analysis}
\label{sec:setup:stats}

The analysis plan was fixed before the comparisons were computed and is reported as
pre-registered. The primary test is a seed-matched paired $t$-test on the differences,
with normality checked by the Shapiro--Wilk test; where normality is rejected, the Wilcoxon
signed-rank test becomes primary for that comparison and this is stated. The Wilcoxon test
is otherwise reported as a non-parametric confirmation.

The sample is small and the consequences are stated rather than left implicit. With three
seeds the paired $t$-test has two degrees of freedom, the Shapiro--Wilk check has almost no
power to detect non-normality, and the exact Wilcoxon test cannot reach $p < 0.05$ at all.
Three seeds are nevertheless informative here because the quantity being tested has a
standard deviation of $0.0015$ while the differences of interest are ten to twenty times
that; the bootstrap intervals reported alongside every comparison are the more robust
summary and should be read as the primary evidence of magnitude. Only the baseline was
repeated over ten seeds, and no claim rests on a three-seed difference that the bootstrap
interval does not separate from zero. Multiplicity is controlled by the
Holm procedure over the four pre-registered comparisons.

Every comparison reports the mean difference, a bootstrap 95\% confidence interval from
10{,}000 resamples, and effect sizes (paired Cohen's $d$ and Cliff's $\delta$) alongside
the $p$-value. Because seed-to-seed variance is small, statistical detectability alone is
a weak criterion: with a baseline standard deviation of $0.0015$ mAP$_{50\text{-}95}$ over
ten seeds, differences an order of magnitude smaller than any reported architectural gain
are still detectable. A floor of $0.005$, roughly three baseline standard deviations, was
therefore fixed in advance: a difference below it is reported as detectable but practically
negligible, and no claim rests on such a difference. The threshold was chosen from the
measured variance of the baseline before any comparison was computed, not from the
comparisons themselves.

\subsection{Data-centric control study}
\label{sec:setup:enrichment}

Before the cost study, a separate controlled experiment tested whether enriching the
training data helps. Starting from the same baseline checkpoint of the same seed, five
conditions were trained for an identical budget of 20 additional epochs, differing only in
the training data: unchanged data; self-training pseudo-labels; pasting of randomly chosen
gallery crops; pasting of crops retrieved by nearest-neighbour search in CLIP embedding
space; and the complete enrichment pipeline. Ten seeds were run per condition, for 60 runs
in total. The random-paste condition shares the placement engine, budget, scale rules and
filters with the retrieval conditions and differs only in how the neighbour is chosen,
which is what makes it a control for compositing artefacts rather than for semantics.
A budget sweep over the number of pasted instances per image and a from-scratch replication
were run as follow-ups. Section~\ref{sec:results} reports the outcome.

\subsection{Implementation and environment}
\label{sec:setup:env}

All experiments run on a single NVIDIA RTX 4500 Ada Generation (24\,GB) under WSL2 with
Ubuntu, Python 3.14, PyTorch 2.13 with CUDA 13.0, and Ultralytics 8.4.118. Library versions
are recorded with the results.

Every reported number is produced by a script and stored in a file; no value in this paper
was transcribed by hand. One row per run records the condition, seed, resolution,
per-image batch, accuracy metrics, measured latency and power, and the number of epochs
actually trained. Figure sources, per-class results, calibration curves and the
enrichment audit are written to separate files at the time each run completes, so that
every figure can be redrawn without repeating a single GPU hour.
